\section{What a transposition construction would need}\label{sec:construction}

The native examples distinguish inexpensive execution from inexpensive transposition.
Expansion of the incidence graph and connected demand pairs can occur in linear time,
but neither feature alone defeats a storage-code bound. Conversely, a construction
cannot be excluded by a complex-linear theorem unless its actual genie fibres lie in
the theorem's linear class with the stated coefficient and block-width bounds.

A candidate oblivious linear-time transposer must therefore fit both sides of the
interface: its trace must evade the available cheap depth-reduction theorems, or its
full finite-alphabet recovery relation must support enough codewords to absorb the
cost of every applicable genie. The connected-demand example shows why a positive
coding deficit uniform in fixed depth is useful for this comparison. It does not
reduce an arbitrary candidate to the row interfaces analyzed here.

For an adaptive construction, fixing one trace is insufficient: any comparison must
also account for the information needed to specify that trace and its recovery sets.
The following section states these remaining obligations precisely.
